\documentclass{scrartcl}
\usepackage[utf8]{inputenc}
\usepackage[T1]{fontenc}
\usepackage{geometry}
\usepackage{graphicx}
\usepackage{float}
\usepackage{hyperref}
\usepackage[english,indonesian]{babel}

\title{Optimasi Pemrosesan Template LaTeX Menggunakan Pendekatan Berbasis AI}
\author{Academic Support Agent}
\date{\today}

\begin{document}
\maketitle

\begin{abstract}

\end{abstract}

\chapter{PENDAHULUAN}


\chapter{TINJAUAN PUSTAKA}


% Original content (commented out):
% \documentclass[a4paper,12pt,oneside]{scrartcl} 

\usepackage[german,ngerman]{babel}
 %\usepackage[utf8]{inputenc}
\usepackage{fontspec} 
 
%  \usepackage{palatino}  %Schriftart schönes als die Übliche
% \usepackage{mathpazo}  %Palatino-Schriftart in Formeln
% \usepackage{microtype} % Trennstriche auf Blocksatzende

\setmainfont{texgyrepagella}[
  Extension=.otf,
  UprightFont=*-regular,
  ItalicFont=*-italic,
  BoldFont=*-bold,
  BoldItalicFont=*-bolditalic,
  SlantedFont=*-regular,
  SlantedFeatures={FakeSlant=0.15},
]
\setsansfont{TeX Gyre Heros}
\setmathrm[Ligatures=TeX]{TeX Gyre Pagella}
\usepackage{unicode-math}
\setmathfont[math-style=ISO,bold-style=ISO,]{TeX Gyre Pagella Math}

\input{preambel_spr} 	% Alle Packages und Formatierungen werden in preamble.tex eingetragen um den Text hier übersichtlicher zu gestalten.
 


\begin{document}
 
%\input{title}			%Schöneres Deckblatt ohne Note
\input{title2}			%leicht modifiziertes offizielles Deckblatt vom KM
 
\begingroup
\pagestyle{empty} % ohne Seitenzahl
\tableofcontents % Inhaltsverzeichnis
\clearpage % Seite bis zum Ende leeren
\endgroup

\newpage

\setcounter{page}{1} % Seitenzähler wird auf 1 gesetzt

\part{Einleitung} % Überschrift für erstes Kapitel, optional auch: [Eintrag im IHV]{Einleitung}
\label{sec:Einleitung} % Textmarke, um darauf zu verweisen.
\lipsum[2]			% Dummytext (einfach löschen und los schreiben)
Hier\autocite[123]{biblatex} sollte die Literatur von Rühmann\autocite{buch1} erwähnt werden, mal sehen, dass die Linie und das Verzeichnis richtig gestellt werden.\footcite[44]{buch1}

\zitatblock{Ein langes Zitat sollte eigentlich eingerückt werden. [\dots] Normalerweise sind ein paar Zeilen kein Problem, aber wenn es an die fünf Zeilen lang ist, dann fällt es durch die Einrückung deutlicher auf. Normalerweise sind ein paar Zeilen kein Problem, aber wenn es an die fünf Zeilen lang ist, dann fällt es durch die Einrückung deutlicher auf. [\dots] Normalerweise sind ein paar Zeilen kein Problem, aber wenn es an die fünf Zeilen lang ist, dann fällt es durch die Einrückung deutlicher auf.\autocite[123]{zitat}}

Hier wird auf das Zitat eingegangen. Dennoch kann hier noch ein \zitat{Zitat aus der Literatur} stehen. Es kann einfach in den Text eingebunden werden.

\part{Thematik in der Literatur}
Wird allerdings komplett wörtlich übernommen so kann \enquote{die heilige Schrift von damals}\autocite{bibel} in Anfühungszeichen übernommen werden
\lipsum[30]				% Dummytext (einfach löschen und los schreiben)

\begin{figure}[ht]		% Bild mit Unterschrift einbinden
\centering
\includegraphics[scale=1]{pics/bild.png} 
\caption[Demobild]{Einzelne Bilder wie dieses sollten zentriert auf der Seite platziert werden. Auf eine sinnvolle Größe ist ebenso zu achten, wie auf eine konkrete Beschreibung der Abbildung.}
\label{fig:bild1}\vspace{0.8em}
\end{figure}

\lipsum[31]

\newpage
\section{Überschrift}

\renewcommand{\baselinestretch}{1.0}\normalsize 
\setlength{\columnsep}{30pt} % Abstand Spalten vergrößern
\begin{multicols}{2}[\noindent \textbf{Hier könnte beispielsweise  ein Verweis auf den Textauszug aus der Literatur stehen.}] %\textbf{Dies kann als Zusammenfassung für den darunter liegenden Text gesehen werden. Hin und wieder ist dieser auch fett gedruckt.}
\setcounter{linenumber}{1}
\begin{linenumbers} % Schaltet Zeilennummerierung ein
%\modulolinenumbers[5] % nur jede 5. Zeile
\lipsum[99]
\end{linenumbers}
\end{multicols} % Ende 2 Spalten
\renewcommand{\baselinestretch}{1.5}\normalsize 

\subsection{Überschrift}
\lipsum[99]

\begin{table}[!ht]
\centering
\caption[Zahlenvergleich der Werte 1 und 2]{Die Werte sind an fiktiven Proben gemessen.}\vspace{0.8em}
\label{tab:tab1}
\begin{tabular}{l|cc} \toprule
 &\multicolumn{2}{r}{universal}\\ \cmidrule{1-1}
\backslashbox{ABC}{DEF} & 1 & 2 \\ \midrule 
E & 3,34 & 4,34 \\
F & 5,45 & 6,45 \\ \bottomrule
\end{tabular}
\end{table}

\lipsum[99]

\subsubsection{Überschriftlein} % tiefere Unterteilung nicht sinnvoll!!!

\newpage
\part{Ausblick und Zusammenfassung}
Dies ist eine Text\footnote{Dies solle Fußnote werden, die den Rahmen einer Zeile sprengt. Dies lässt sich auch durch Wiederholung den Inhalts erreichen. Dies lässt sich auch durch Wiederholung den Inhalts erreichen.  Dies lässt sich auch durch Wiederholung den Inhalts erreichen. \newline Dies lässt sich auch durch Wiederholung den Inhalts erreichen.}, der auch in der Literatur existiert.\autocite[23]{demo} 
\lipsum[32]

Da pack ich doch gleich no ne Fußnote drunter.\footnote{Mal sehen, welcher Abstand hier zu vorherigen Fußnote entsteht.}

\lipsum[10]

\begin{table}[!ht]					%	Tabelle genau hier einsetzen
\centering							% Tabelle zentrieren
\caption[YYY]{Tabellenüberschrift.}\vspace{0.8em} % Tabellenbeschrift.
\label{tab:2}							% Label für den Tabellenverweis
\begin{tabular}{l|r|r} \toprule		% l = linksb., c = zentriert, r = rechtsb., | vert. Strich
  & \textbf{Spalte A} & \textbf{Spalte B} \\ \midrule 				% & ist Spaltentrennzeichen, Abschluss der Zeile mit "\\"
\textbf{Zeile 1} & Wert 1A & Wert 1B \\
\textbf{Zeile 2 }& Wert 2A & Wert 2B \\ \bottomrule		% top-, mid- & bottomrule in untersch. Strichstärken
\end{tabular}
\end{table}

\lipsum[41]

\input{selbst}      % selbst.tex wird hier integriert = Selbstständigkeitserklärung... bitte noch nach eigenen Vorstellungen abändern

\newpage
\listoffigures		% Abbildungsverzeichnis erstellen

\newpage
\listoftables		% Tabellenverzeichnis erstellen

\newpage
\addcontentsline{toc}{section}{Literaturverzeichnis}
\small 						%\documentclass[journal = jacsat]{achemso}
							%\setkeys{acs}{articletitle = true}
%\printbibliography
\renewcommand{\baselinestretch}{1.0}\normalsize
\printbibheading[title={Literaturverzeichnis}]%[title=Blub]
\printbibliography[heading=subbibliography,type=book, title={Bücher}]
\printbibliography[heading=subbibliography,type=article, title={Artikel}]
%\printbibliography[heading=subbibliography,type=online, title={Internetquellen}]
\printbibliography[heading=subbibliography,type=manual, title={Anleitungen}]
\printbibliography[heading=subbibliography,type=report, title={Reports}]
 
\newpage
\appendix
\addcontentsline{toc}{section}{Anhang}
\begin{appendix}
\section*{Anhang}
%\subsection{Anhangsteil 1}
%\lipsum[31]
%\subsection{Anhangsteil 2}
\end{appendix}


 
\end{document}

\end{document}